\chapter{準備}
\label{ch:preliminaries}

本章では、以降の章で極大クリーク列挙アルゴリズムを並列化するために必要な概念を定義する。
グラフとクリークの基本用語から始め、逐次アルゴリズムである Bron--Kerbosch 法とそのピボット版 CLIQUES、疎グラフ向けの Eppstein アルゴリズムを説明する。
続いて実験に用いる十種類のグラフを一覧で示し、最後に第\ref{ch:implementation}章以降で使う Python の multiprocessing の前提だけを述べる。

\section{グラフとクリーク}
\label{sec:graph-clique}

本論文では、無向グラフ $G = (V, E)$ を扱う。
$V$ は頂点の集合、$E \subseteq \{\, \{u, v\} \mid u, v \in V,\ u \neq v \,\}$ は辺の集合であり、頂点数を $n = |V|$、辺数を $m = |E|$ と書く。
自己ループと多重辺は持たないものとする。
頂点 $v$ の近傍を $N(v) = \{\, u \in V \mid \{u, v\} \in E \,\}$ で表し、その大きさ $|N(v)|$ を $v$ の次数と呼ぶ。

頂点部分集合 $C \subseteq V$ が\textbf{クリーク}であるとは、$C$ の任意の二頂点が辺で結ばれていること、すなわち $C$ の誘導部分グラフが完全グラフであることをいう。
クリーク $C$ が\textbf{極大クリーク}であるとは、$C$ にどの頂点を加えてもクリークでなくなること、言い換えれば $C \subsetneq C'$ となるクリーク $C'$ が存在しないことをいう。
極大クリークは、頂点数が最大のクリークを指す\textbf{最大クリーク}とは異なる概念である。
最大クリークは常に極大クリークだが、その逆は成り立たない。
本論文が対象とするのは、グラフ中のすべての極大クリークを重複なく出力する\textbf{極大クリーク列挙}である。

\section{Bron--Kerbosch 法とピボット版 CLIQUES}
\label{sec:bron-kerbosch}

極大クリーク列挙の基礎となるのが Bron--Kerbosch 法である。
この方法は、三つの頂点集合を持って再帰的に探索する。
確定済みの現在のクリーク $R$、まだ $R$ に加えられる候補 $P$、すでに調べ終えて再訪を禁じる $X$ である。
$P$ と $X$ がともに空になったとき、$R$ はそれ以上広げられず他の枝でも数えられない極大クリークなので、これを出力する。
候補 $P$ の各頂点 $v$ について、$v$ を $R$ に加え、$P$ と $X$ を $v$ の近傍だけに絞って再帰し、戻ったら $v$ を $P$ から $X$ へ移す。

素朴な Bron--Kerbosch 法は、同じ極大クリークに至る枝を何度も展開してしまう。
この無駄を抑えるのが\textbf{ピボット}の選択である。
候補と処理済みの合併 $P \cup X$ から一頂点 $u$（ピボット）を選ぶと、任意の極大クリークは $u$ かその非近傍のいずれかを必ず含むため、再帰を $P$ 全体ではなく「$u$ の非近傍」だけに限定できる。
Tomita, Tanaka, Takahashi は、$|P \cap N(u)|$ を最大にする $u$ をピボットに選ぶ版を\textbf{CLIQUES}として定式化した\cite{tomita2006}。
この選び方は再帰の分岐を最も強く刈り取り、最悪計算量を $O(3^{n/3})$ に抑える。

この計算量は、出力しうる極大クリークの個数という下限に照らして最善である。
$n$ 頂点の無向グラフが持ちうる極大クリークの個数は高々 $3^{n/3}$ であり、この上限は頂点を3個ずつの組に分けた完全多部グラフ $K_{3,3,\dots,3}$ が達成することを Moon と Moser が示した\cite{moon1965}。
極大クリークを一つ出力するのに最低でも定数時間かかる以上、列挙にかかる時間はこの個数を下回れない。
CLIQUES の $O(3^{n/3})$ はこの上限と一致しており、最悪の入力に対して漸近的に最適である\footnote{全体計算量と、極大クリークを一つ出すごとの遅延（delay）の関係は Conte と Tomita による後続の解析で詳しく調べられている\cite{conte2022}。本論文の逐次実装はこの論文の擬似コードに従う。}。

\section{縮退数と縮退順序}
\label{sec:degeneracy}

CLIQUES の最悪計算量は頂点数 $n$ だけで決まり、グラフが疎であることを利用しない。
現実の多くのグラフは辺が少なく、この疎さを計算量に反映できれば列挙は速くなる。
その疎さを測る尺度が\textbf{縮退数}である。

グラフ $G$ が $d$-縮退であるとは、$G$ のどの部分グラフにも次数 $d$ 以下の頂点が存在することをいう。
これが成り立つ最小の $d$ を $G$ の縮退数と呼ぶ。
縮退数は\textbf{縮退順序}を通しても特徴づけられる。
次数が最小の頂点を繰り返し取り除き、取り除いた順に並べたものが縮退順序であり、この順序では、どの頂点も自分より後ろに位置する近傍（後方近傍）を高々 $d$ 個しか持たない。
縮退順序は最小次数の頂点を順に剥がすだけで $O(n + m)$ 時間で構成でき、疎なグラフほど $d$ が小さい。

\section{Eppstein アルゴリズム}
\label{sec:eppstein}

Eppstein, L\"offler, Strash は、縮退順序を使って CLIQUES の計算量を縮退数で表す変種を与えた\cite{eppstein2010}。
最外層の再帰を縮退順序に沿って次数の低い頂点から回し、各頂点ではその後方近傍だけを候補に取る。
後方近傍は高々 $d$ 個なので、各頂点を根とする部分問題は最大 $d$ 頂点のグラフ上の列挙に帰着し、その内部ではピボット選択を Tomita 流のまま用いる。
これにより全体の計算量は $O(d \cdot n \cdot 3^{d/3})$ となり、$d$ が小さい疎グラフでは $n$ にほぼ比例する時間まで下がる。
実データを含む大規模な疎グラフでの有効性は、続く実験版の論文で確かめられている\cite{eppstein2013}。

このアルゴリズムの構造は、本論文の並列化にとって重要である。
最外層が縮退順序の各頂点を根とする $n$ 個の独立した部分問題に分かれるため、部分問題を worker へ配るだけで最外層を並列化できる。
$d$ が小さいほど各部分問題は小さく独立性が高い。

\section{実験に用いるグラフ}
\label{sec:graphs}

実験では、次数分布と疎密の異なる十種類の無向グラフを用いる。
生成規則が既知の合成グラフ七種と、実データである SNAP データセット三種に分かれる。
比較の主軸は縮退数 $d$ であり、$d$ が小さいグラフほど疎で、Eppstein アルゴリズムが最外層を小さく独立した部分問題に分割しやすい。

合成グラフ七種は、疎密に加えて次数分布と出力サイズという二つの軸でも選んでいる。
次数分布については、次数がほぼ均一な Erd\H{o}s--R\'enyi と、少数の頂点に次数が集中する Barab\'asi--Albert（ソーシャルネットワークに近いべき乗則分布）を対比させる。
出力サイズについては、辺の数だけ極大クリークが出る2次元格子と、理論上限 $3^{k}$ 個をそのまま達成する Moon--Moser を両端に置く。
各グラフの内容は次のとおりである。

\begin{itemize}
  \item \textbf{木}：頂点 $i$ が、それより前の頂点から一様ランダムに一つ選んで接続する再帰的ランダム木。三角形を持たず、$d = 1$。
  \item \textbf{2次元格子}：$100 \times 100$ の4近傍格子。三角形を持たないため、辺の一本ずつがそのまま極大クリークになる。$d = 2$。
  \item \textbf{疎 Erd\H{o}s--R\'enyi}：平均次数を10に固定して頂点対を一様ランダムに結ぶ $G(n, m)$ 型のグラフ。$d = 7$。
  \item \textbf{スケールフリー Barab\'asi--Albert}：新規頂点が既存頂点の次数に比例した確率で $m = 5$ 本の辺を張る優先的接続モデル。$d = 5$。
  \item \textbf{密 Erd\H{o}s--R\'enyi}：頂点対ごとに確率 $p = 0.5$ で辺を張る $G(n, p)$ モデル。辺数が $n^2$ で効くため小さい $n$ でのみ使う。$d = 48$。
  \item \textbf{Moon--Moser}：頂点を3個ずつ $k = 11$ 組に振り分けた完全多部グラフ $K_{3,3,\dots,3}$。頂点数 $3k$ に対し極大クリークが $3^{k}$ 個生じる、\ref{sec:bron-kerbosch}節で述べた極値グラフである。$d = 30$。
  \item \textbf{完全グラフ}：全頂点対に辺がある $K_{300}$。極大クリークは全体で1個だけになる。$d = 299$。
\end{itemize}

実データは、Stanford の SNAP（Stanford Large Network Dataset Collection）が配布する無向グラフ三種である。
いずれも頂点対の行を並べたエッジリスト形式で読み込み、有向や符号つきの情報は捨てて無向グラフとして扱う。
ca-AstroPh と ca-HepPh は arXiv の共著関係のネットワークで、それぞれ天体物理学分野と高エネルギー物理学分野に対応する。
soc-sign-epinions は商品レビューサイト Epinions のユーザー間の信頼・不信ネットワークで、符号を無視して無向グラフとして扱う。
このなかで soc-sign-epinions は頂点数と極大クリーク数がともに大きく、逐次で約95秒かかるため、第\ref{ch:experiments}章では並列化の効果を測る主力データとして用いる。

各グラフの頂点数 $n$、縮退数 $d$、極大クリーク数を表\ref{tab:graphs}に示す。

\begin{table}[htbp]
  \centering
  \caption{実験に用いるグラフ。上段が合成グラフ七種、下段が SNAP データセット三種。}
  \label{tab:graphs}
  \begin{tabular}{lrrr}
    \toprule
    グラフ & $n$ & $d$ & 極大クリーク数 \\
    \midrule
    木 & 20{,}000 & 1 & 20.0千 \\
    2次元格子 $100 \times 100$ & 10{,}000 & 2 & 19.8千 \\
    疎 Erd\H{o}s--R\'enyi（平均次数10） & 20{,}000 & 7 & 99.7千 \\
    スケールフリー Barab\'asi--Albert（$m=5$） & 20{,}000 & 5 & 96.9千 \\
    密 Erd\H{o}s--R\'enyi（$n=120$, $p=0.5$） & 120 & 48 & 3.5万 \\
    Moon--Moser（$k=11$） & 33 & 30 & 17.7万 \\
    完全グラフ $K_{300}$ & 300 & 299 & 1 \\
    \midrule
    ca-AstroPh & 18{,}772 & 56 & 36.4千 \\
    ca-HepPh & 12{,}008 & 238 & 14.9千 \\
    soc-sign-epinions & 131{,}828 & 121 & 2{,}223万 \\
    \bottomrule
  \end{tabular}
\end{table}

\todo{合成グラフ七種の外観図（Notion に描画済み画像あり）を挿入する。}

\section{Python multiprocessing による並列化}
\label{sec:multiprocessing}

第\ref{ch:implementation}章以降の並列化は Python の標準ライブラリ multiprocessing を用いる。
CPython はインタプリタ全体を一つのロック（GIL）で保護しており、複数スレッドが Python バイトコードを同時に実行できない。
そのため CPU バウンドな列挙を並列化するには、スレッドではなく複数のプロセスを使い、各プロセスに独立したインタプリタを持たせる必要がある。
以下では、並列化のコストを左右するプロセス生成方式の違いだけを述べる。

子プロセスの生成方式には \texttt{spawn} と \texttt{fork} がある。
\texttt{spawn} は子プロセスで新しいインタプリタを起動し、worker に渡す引数を親側で pickle により直列化して送り、子側で復元する。
\texttt{fork} は親プロセスのメモリを copy-on-write で子に継承させるため、引数の直列化も転送も起こらない。
macOS の既定は \texttt{spawn} であり、これは worker へ渡すグラフを worker 数だけ pickle 転送することを意味する。
このグラフ転送が並列化の固定費として効いてくる点は、第\ref{ch:experiments}章で実測とともに扱う。
